\scanpage{68}
[Cela est le cas en particulier dans le cas de l'exemple I, \deleted{$S'$ étant} \inserted{lorsque $S'$ est} fini et loc. libre sur $S$, et l'exemple II, \inserted{lorsque} $W$ \deleted{\unsure{de|est}} est affine sur $S$].

\uline{Proposition} Sous ces conditions, si \deleted{$\unsure{\mathfrak X|\mathcal X}$ est} $\mathfrak X$ est un schéma \uline{affine} sur $\underline{\mathbb W}(S)$, $\mathbb W^a(\mathfrak X)$ existe et est un schéma \uline{affine} sur $S$.

\uline{Démonstration} \deleted{Soit $\mathfrak X$} \inserted{On peut supposer $S$ affine.} Si $\mathfrak X=\operatorname{Spec}(\unsure{\mathcal A|\mathfrak A|\mathcal B|\mathcal C})$ où $\mathcal A$ est une algèbre de type fini \deleted{\unsure{de}} sur $\underline{\mathbb W}(A)$, engendrée par $r$ éléments, il résulte des \unsure{sorites|résultats} du \S{} N\textsuperscript{o} 1 que $\mathbb W^a(\mathfrak X)$ existe et est $\simeq$ un sous-préschéma fermé de $\mathbb W^r$ [défini par des ``équations \unsure{polynômiales|polynôminales|polynomiales} \unsure{à|:}\inserted{\unsure{de|dans}} coeff. dans $\Gamma(\mathbb W/S)$]. Comme $\mathbb W$ est affine, il en est de $\unsure{=|\simeq}$ de \deleted{$\mathbb W$} $\mathbb W^r$ donc aussi de $\mathbb W^a(\mathfrak X)$ (\unsure{ouf|cqf|vif|cqfd}). On a obtenu d'ailleurs

\uline{Cor}. \deleted{Si} \deleted{\unsure{Alors} $\mathcal A$ \unsure{faisceau de}} \inserted{Supposons $S$ affine.} Si $\mathfrak X$ est engendré par $r$ générateurs sur celui de $\mathbb W(S)$, et celui de $\mathbb W$ par $s$ gén. sur celui de $S$, alors
% archive sheet number: 67
